% !Mode:: "TeX:UTF-8"
% 任务书中的信息
%% 原始资料及设计要求
\assignReq
{}
{}
{}
{}
{}
%% 工作内容
\assignWork
{本文系统介绍了研究Calabi‑Yau簇形变理论的两个主要途径。复几何方法上，}
{我们从 Maurer–Cartan 方程出发，建立了 Kodaira–Spencer 映射，并证明光滑 }
{Calabi‑Yau流形的形变无障碍。
代数方法上，我们构造了 \(T^i\) 函子与余切复形，
}
{
给出了概形形变的障碍理论，用Schlessinger判别法证明了形变函子一定条}
{件下的pro-可表示性，
并在此基础上介绍了\(T^1\) 提升性质及其在复射影光滑 }
{Calabi‑Yau 簇上的应用。}
%% 参考文献
\assignRef
{R. Hartshorne. Algebraic Geometry, volume 52 of Graduate Texts in Mathematics.}
{R. Hartshorne. Deformation Theory, volume 257 of Graduate Texts in Mathematics.}
{S. Lichtenbaum and M. Schlessinger. The cotangent complex of a morphism.}
{R. Friedman. Simultaneous resolution of threefold double points. Math. Ann.}
{Z. Ran. Deformations of Calabi–Yau threefolds with ordinary double points.}
{}
{}
{}
